% ================================================================
%  SESIÓN 09 — Perímetro y área: no son lo mismo
%  Almería en Cifras y Formas · Matemáticas 1.º ESO
%  IES Al-Ándalus · 2025-2026
% ================================================================
\documentclass[aspectratio=169, 10pt, spanish]{beamer}
\usepackage{almeria-beamer}

\renewcommand{\SessionNumber}{09}
\renewcommand{\SessionTitle}{Perímetro y área: no son lo mismo}
\renewcommand{\SessionName}{Sesión 09}
\renewcommand{\SessionObjective}{%
  Distinguir conceptual y numéricamente el perímetro del área,
  comprendiendo que son medidas independientes y que una puede
  cambiar sin afectar a la otra.%
}
\renewcommand{\BlockName}{Bloque 2: Geometría urbana}

\begin{document}

% ---------------------------------------------------------------
%  PORTADA
% ---------------------------------------------------------------
\begin{frame}[plain]
  \titlepage
\end{frame}

% ---------------------------------------------------------------
%  ÍNDICE
% ---------------------------------------------------------------
\begin{frame}{Contenidos de hoy}
  \begin{columns}[t]
    \begin{column}{0.48\textwidth}
      \begin{itemize}
        \item Tabla comparativa P vs.\ A
        \item Mismos P, distintos A
        \item Mismos A, distintos P
        \item Problema isoperimérico
        \item Implicaciones en urbanismo
      \end{itemize}
    \end{column}
    \begin{column}{0.48\textwidth}
      \begin{itemize}
        \item Almería: Rambla completa
        \item Cálculo paralelo P y A
        \item Ejercicios Bloom
        \item Error frecuente
        \item Evidencia del día
      \end{itemize}
    \end{column}
  \end{columns}
  \vspace{0.4cm}
  \begin{notabox}[Niveles de Bloom]
    \bloomtag{Recordar}\quad
    \bloomtag{Comprender}\quad
    \bloomtag{Aplicar}
  \end{notabox}
\end{frame}

% ---------------------------------------------------------------
\seccionframe{Comparando perímetro y área}
% ---------------------------------------------------------------

% ---------------------------------------------------------------
%  FRAME 3 — Tabla comparativa
% ---------------------------------------------------------------
\begin{frame}{Perímetro vs. Área: tabla comparativa}
  \begin{center}
  \AlmFontSmall
  \begin{tabular}{|>{\bfseries\color{AlmAzulM}}l|>{\centering\arraybackslash}p{3.8cm}|>{\centering\arraybackslash}p{3.8cm}|}
    \hline
    \rowcolor{AlmAzul!20}
    & \textbf{\color{AlmTerra}Perímetro} & \textbf{\color{AlmVerde}Área} \\
    \hline
    ¿Qué mide? &
      El \textbf{contorno} &
      La \textbf{superficie}\\ &
      (la valla) &(el césped) \\
    \hline
    Dimensión &
      \textbf{1D} (lineal) &
      \textbf{2D} (superficial) \\
    \hline
    Unidades &
      m, km, cm\ldots &
      m$^2$, km$^2$, ha\ldots \\
    \hline
    Fórmula rect. &
      $P = 2(l + w)$ &
      $A = l \times w$ \\
    \hline
    Uso &
      Vallas, bordes, caminos &
      Suelos, pinturas, césped \\
    \hline
  \end{tabular}
  \end{center}

  \vspace{0.2cm}
  \begin{notabox}[Idea clave]
    \AlmFontSmall Perímetro y área miden cosas \textbf{completamente distintas}.
    Una puede ser grande mientras la otra es pequeña.
  \end{notabox}
\end{frame}

% ---------------------------------------------------------------
%  FRAME 4 — Mismo P, distinto A
% ---------------------------------------------------------------
\begin{frame}{Mismo perímetro, distinto área}
  \begin{columns}[c]
    \begin{column}{0.62\textwidth}
      \begin{center}
      \begin{tikzpicture}[scale=0.78]
        % Rectángulo 6x2
        \fill[AlmTerra!12] (0,0) rectangle (4.8,1.6);
        \draw[rectot, line width=1.8pt] (0,0) rectangle (4.8,1.6);
        \draw[acota] (0,-0.55) -- (4.8,-0.55)
          node[midway,below,font=\AlmFontScript]{6};
        \draw[acota] (5.15,0) -- (5.15,1.6)
          node[midway,right,font=\AlmFontScript]{2};
        \node[font=\AlmFontScript\bfseries, color=AlmTerra] at (2.4,0.8)
          {$P = 16$ \quad $A = 12$};

        % Rectángulo 4x4 (más abajo)
        \fill[AlmVerde!15] (0,-3.0) rectangle (3.2,-3.0+3.2);
        \draw[rectov, line width=1.8pt] (0,-3.0) rectangle (3.2,-3.0+3.2);
        \draw[acota] (0,-3.6) -- (3.2,-3.6)
          node[midway,below,font=\AlmFontScript]{4};
        \draw[acota] (3.55,-3.0) -- (3.55,-3.0+3.2)
          node[midway,right,font=\AlmFontScript]{4};
        \node[font=\AlmFontScript\bfseries, color=AlmVerde] at (1.6,-1.4)
          {$P = 16$ \quad $A = 16$};

        % Etiqueta
        \node[font=\AlmFontScript\bfseries, color=AlmAmbar!80!black]
          at (2.5,-0.1) {};
      \end{tikzpicture}
      \end{center}
    \end{column}
    \begin{column}{0.34\textwidth}
      \begin{teoriabox}
        \AlmFontSmall
        Ambos rectángulos tienen
        $P = 16$ unidades.
        \\[6pt]
        Pero sus áreas son distintas:
        \begin{itemize}
          \item 6$\times$2: $A = 12$
          \item 4$\times$4: $A = 16$
        \end{itemize}
        \vspace{4pt}
        El cuadrado es la figura
        de \textbf{mayor área} para
        un perímetro fijo.
        \\[4pt]
        \AlmFontScript(Problema isoperimérico)
      \end{teoriabox}
    \end{column}
  \end{columns}
\end{frame}

% ---------------------------------------------------------------
%  FRAME 5 — Mismo A, distinto P
% ---------------------------------------------------------------
\begin{frame}{Mismo área, distinto perímetro}
  \begin{columns}[c]
    \begin{column}{0.62\textwidth}
      \begin{center}
      \begin{tikzpicture}[scale=0.78]
        % Rectángulo 3x4 → A=12, P=14
        \fill[AlmAzulC!18] (0,0) rectangle (2.4,3.2);
        \draw[zona, line width=1.8pt] (0,0) rectangle (2.4,3.2);
        \draw[acota] (0,-0.55) -- (2.4,-0.55)
          node[midway,below,font=\AlmFontScript]{3};
        \draw[acota] (2.75,0) -- (2.75,3.2)
          node[midway,right,font=\AlmFontScript]{4};
        \node[font=\AlmFontScript\bfseries, color=AlmAzulM] at (1.2,1.6)
          {$P = 14$\quad$A = 12$};

        % Rectángulo 6x2 → A=12, P=16
        \fill[AlmAmbar!15] (4.2,0.8) rectangle (4.2+4.8,0.8+1.6);
        \draw[AlmAmbar!70!black, line width=1.8pt]
          (4.2,0.8) rectangle (4.2+4.8,0.8+1.6);
        \draw[acota] (4.2,0.25) -- (9.0,0.25)
          node[midway,below,font=\AlmFontScript]{6};
        \draw[acota] (9.35,0.8) -- (9.35,2.4)
          node[midway,right,font=\AlmFontScript]{2};
        \node[font=\AlmFontScript\bfseries, color=AlmAmbar!80!black] at (6.6,1.6)
          {$P = 16$\quad$A = 12$};
      \end{tikzpicture}
      \end{center}
    \end{column}
    \begin{column}{0.34\textwidth}
      \begin{teoriabox}
        \AlmFontSmall
        Ambos rectángulos tienen
        $A = 12$ unidades cuadradas.
        \\[6pt]
        Pero sus perímetros son distintos:
        \begin{itemize}
          \item 3$\times$4: $P = 14$
          \item 6$\times$2: $P = 16$
        \end{itemize}
        \vspace{4pt}
        La figura más «cuadrada»
        tiene menor perímetro para
        una misma área.
      \end{teoriabox}
    \end{column}
  \end{columns}
\end{frame}

% ---------------------------------------------------------------
%  FRAME 6 — Implicaciones en urbanismo
% ---------------------------------------------------------------
\begin{frame}{¿Cuándo usar cada medida?}
  \begin{columns}[t]
    \begin{column}{0.48\textwidth}
      \begin{definicionbox}{Necesitas el PERÍMETRO}
        \AlmFontSmall
        \begin{itemize}
          \item Colocar una \textbf{valla} alrededor de un parque.
          \item Poner \textbf{bordillo} en el borde de una acera.
          \item Calcular la distancia de una \textbf{ruta a pie}.
          \item Instalar \textbf{luminarias} en el contorno de una plaza.
        \end{itemize}
      \end{definicionbox}
    \end{column}
    \begin{column}{0.48\textwidth}
      \begin{definicionbox}{Necesitas el ÁREA}
        \AlmFontSmall
        \begin{itemize}
          \item Comprar \textbf{baldosas} para pavimentar.
          \item Calcular el \textbf{césped} a sembrar.
          \item Valorar el \textbf{precio} de un terreno.
          \item Pintar \textbf{la fachada} de un edificio.
        \end{itemize}
      \end{definicionbox}
    \end{column}
  \end{columns}
  \vspace{0.3cm}
  \begin{notabox}[Resumen]
    \AlmFontSmall \textbf{Valla, bordillo, camino} $\rightarrow$ Perímetro (metros).
    \quad
    \textbf{Suelo, césped, pintura} $\rightarrow$ Área (m$^2$).
  \end{notabox}
\end{frame}

% ---------------------------------------------------------------
\seccionframe[BAplicar]{Almería: la Rambla}
% ---------------------------------------------------------------

% ---------------------------------------------------------------
%  FRAME 7 — Ejemplo completo: Rambla de Almería
% ---------------------------------------------------------------
\begin{frame}{Ejemplo completo: Rambla de Almería}
  \begin{columns}[c]
    \begin{column}{0.52\textwidth}
      \begin{ejemplobox}
        \AlmFontSmall La \textbf{Rambla de Almería} es aproximadamente
        rectangular con \textbf{370 m} de largo y \textbf{50 m}
        de ancho.
        \begin{align*}
          P &= 2(370 + 50) = 2 \times 420 = \mathbf{840\ m} \\
          A &= 370 \times 50 = \mathbf{18{,}500\ m^2}
            = \mathbf{1{,}85\ ha}
        \end{align*}
        \textbf{P} sirve para: vallar todo el contorno. \\
        \textbf{A} sirve para: pavimentar la superficie.
      \end{ejemplobox}
    \end{column}
    \begin{column}{0.44\textwidth}
      \begin{center}
      \begin{tikzpicture}[scale=0.70]
        % Rambla (rectángulo alargado)
        \fill[AlmAmbar!15] (0,0) rectangle (5.5,0.75);
        \draw[AlmAmbar!70!black, line width=2pt] (0,0) rectangle (5.5,0.75);
        % Valla (borde rojo)
        \draw[AlmTerra, line width=3pt, dashed] (0,0) rectangle (5.5,0.75);
        % Árboles (zona verde interior simplificada)
        \fill[AlmVerde!20] (0.2,0.1) rectangle (5.3,0.65);
        \draw[acota] (0,-0.5) -- (5.5,-0.5)
          node[midway,below,font=\AlmFontTiny]{370 m};
        \draw[acota] (5.9,0) -- (5.9,0.75)
          node[midway,right,font=\AlmFontTiny]{50 m};
        \node[font=\AlmFontTiny\bfseries, color=AlmVerde!70!black] at (2.75,0.375)
          {Rambla de Almería};
        \node[font=\AlmFontTiny\bfseries, color=AlmTerra] at (2.75,-0.95)
          {$P = 840$ m};
        \node[font=\AlmFontTiny\bfseries, color=AlmVerde!70!black] at (2.75,-1.35)
          {$A = 18{.}500$ m$^2$};
      \end{tikzpicture}
      \end{center}
    \end{column}
  \end{columns}
\end{frame}

% ---------------------------------------------------------------
\seccionframe[BRecordar]{Ejercicios — RECORDAR}
% ---------------------------------------------------------------

% ---------------------------------------------------------------
%  FRAME 8 — Ejercicio RECORDAR
% ---------------------------------------------------------------
\begin{frame}{Ejercicio 1 — RECORDAR \dificultad{1}}
  \begin{recordarbox}
    Para un rectángulo de \textbf{12 m} de largo y \textbf{5 m}
    de ancho, calcula:
    \begin{enumerate}
      \item El \textbf{perímetro}. Indica su unidad.
      \item El \textbf{área}. Indica su unidad.
    \end{enumerate}
    ¿Cuál es la diferencia entre las unidades de perímetro
    y área?
  \end{recordarbox}
  \vspace{0.25cm}
  \begin{center}
  \begin{tikzpicture}[scale=0.80]
    \fill[AlmAzulC!12] (0,0) rectangle (4.8,2.0);
    \draw[zona, line width=1.8pt] (0,0) rectangle (4.8,2.0);
    \draw[acota] (0,-0.55) -- (4.8,-0.55)
      node[midway,below,font=\AlmFontScript]{12 m};
    \draw[acota] (5.15,0) -- (5.15,2.0)
      node[midway,right,font=\AlmFontScript]{5 m};
    \node[font=\AlmFontSmall, color=AlmTerra] at (1.8,1.0) {$P = ?$\ m};
    \node[font=\AlmFontSmall, color=AlmVerde] at (3.4,1.0) {$A = ?$\ m$^2$};
  \end{tikzpicture}
  \end{center}
\end{frame}

% ---------------------------------------------------------------
\seccionframe[BComprender]{Ejercicios — COMPRENDER}
% ---------------------------------------------------------------

% ---------------------------------------------------------------
%  FRAME 9 — Ejercicio COMPRENDER (problema isoperimérico)
% ---------------------------------------------------------------
\begin{frame}{Ejercicio 2 — COMPRENDER \dificultad{3}}
  \begin{comprenderbox}
    Tienes \textbf{24 m} de valla para cercar un huerto rectangular.
    Explora todas las posibilidades (lados enteros) y encuentra
    la que da \textbf{mayor área}:

    \vspace{0.15cm}
    \AlmFontSmall
    \begin{center}
    \begin{tabular}{|c|c|c|c|}
      \hline
      \rowcolor{AlmAzul!15}
      \textbf{Largo} & \textbf{Ancho} & \textbf{Perímetro} & \textbf{Área} \\
      \hline
      1 m & 11 m & 24 m & \phantom{000} m$^2$ \\
      2 m & 10 m & 24 m & \phantom{000} m$^2$ \\
      3 m &  9 m & 24 m & \phantom{000} m$^2$ \\
      4 m &  8 m & 24 m & \phantom{000} m$^2$ \\
      5 m &  7 m & 24 m & \phantom{000} m$^2$ \\
      6 m &  6 m & 24 m & \phantom{000} m$^2$ \\
      \hline
    \end{tabular}
    \end{center}

    ¿Qué forma geométrica da el área máxima?
  \end{comprenderbox}
\end{frame}

% ---------------------------------------------------------------
\seccionframe[BAplicar]{Ejercicios — APLICAR}
% ---------------------------------------------------------------

% ---------------------------------------------------------------
%  FRAME 10 — Ejercicio APLICAR
% ---------------------------------------------------------------
\begin{frame}{Ejercicio 3 — APLICAR \dificultad{3}}
  \begin{aplicarbox}
    \AlmFontSmall Un urbanista de Almería necesita:
    \begin{enumerate}
      \item \textbf{Vallar} un parque rectangular de 80 m $\times$ 60 m
            (precio: 45 €/m de valla).
      \item \textbf{Cubrir} su superficie con gravilla
            (precio: 12 €/m$^2$).
    \end{enumerate}
    Calcula:
    \begin{enumerate}[a)]
      \item El coste total de la valla.
      \item El coste total de la gravilla.
      \item ¿Cuál es mayor? ¿Qué dato usa cada cálculo (P o A)?
    \end{enumerate}
  \end{aplicarbox}
  \vspace{0.15cm}
  \begin{columns}[c]
    \begin{column}{0.50\textwidth}
      \begin{pasobox}[1]
        \AlmFontScript $P = 2(80 + 60) = ?$ m
        \\
        Coste valla $= P \times 45$~€
      \end{pasobox}
    \end{column}
    \begin{column}{0.46\textwidth}
      \begin{pasobox}[2]
        \AlmFontScript $A = 80 \times 60 = ?$~m$^2$
        \\
        Coste gravilla $= A \times 12$~€
      \end{pasobox}
    \end{column}
  \end{columns}
\end{frame}

% ---------------------------------------------------------------
%  FRAME 11 — Los tres casos TikZ (comparativa completa)
% ---------------------------------------------------------------
\begin{frame}{Tres rectángulos: misma P o misma A}
  \begin{center}
  \begin{tikzpicture}[scale=0.72]
    % --- Caso 1: 6x2, P=16, A=12 ---
    \fill[AlmTerra!12] (0,0) rectangle (4.8,1.6);
    \draw[rectot, line width=1.5pt] (0,0) rectangle (4.8,1.6);
    \node[above, font=\AlmFontTiny\bfseries, color=AlmTerra] at (2.4,1.6)
      {$6 \times 2$};
    \node[font=\AlmFontTiny, color=AlmTerra] at (2.4,0.8) {$P=16$};
    \node[font=\AlmFontTiny, color=AlmTerra] at (2.4,0.3) {$A=12$};

    % --- Caso 2: 4x4, P=16, A=16 ---
    \fill[AlmVerde!15] (5.8,0) rectangle (5.8+3.2,3.2);
    \draw[rectov, line width=1.5pt] (5.8,0) rectangle (5.8+3.2,3.2);
    \node[above, font=\AlmFontTiny\bfseries, color=AlmVerde] at (7.4,3.2)
      {$4 \times 4$};
    \node[font=\AlmFontTiny, color=AlmVerde] at (7.4,1.8) {$P=16$};
    \node[font=\AlmFontTiny, color=AlmVerde] at (7.4,1.2) {$A=16$};

    % Flecha entre 1 y 2
    \draw[->, AlmAmbar!80!black, line width=1.2pt]
      (5.0,0.8) -- (5.6,0.8);
    \node[above, font=\AlmFontTiny\bfseries, color=AlmAmbar!80!black]
      at (5.3,0.8) {P igual};

    % --- Caso 3: 3x4, P=14, A=12 ---
    \fill[AlmAzulC!18] (0,-3.5) rectangle (2.4,-3.5+3.2);
    \draw[zona, line width=1.5pt] (0,-3.5) rectangle (2.4,-3.5+3.2);
    \node[below, font=\AlmFontTiny\bfseries, color=AlmAzulM] at (1.2,-3.5)
      {$3 \times 4$};
    \node[font=\AlmFontTiny, color=AlmAzulM] at (1.2,-2.0) {$P=14$};
    \node[font=\AlmFontTiny, color=AlmAzulM] at (1.2,-2.5) {$A=12$};

    % Flecha entre 3 y 1 (misma A)
    \draw[->, AlmGris, dashed, line width=1.0pt]
      (2.6,-2.1) -- (4.4,-0.5);
    \node[right, font=\AlmFontTiny\bfseries, color=AlmGris]
      at (3.5,-1.4) {A igual};
  \end{tikzpicture}
  \end{center}
  \begin{notabox}[Conclusión]
    \AlmFontSmall Si cambias la forma de un rectángulo manteniendo el mismo
    perímetro, el área cambia. Y viceversa.
  \end{notabox}
\end{frame}

% ---------------------------------------------------------------
%  FRAME 12 — Error frecuente
% ---------------------------------------------------------------
\begin{frame}{¡Cuidado con este error!}
  \begin{errorbox}
    \begin{itemize}
      \item \incorrecto\
            \textbf{«El área del parque es 280 metros»}
            \\[6pt]
            \correcto\
            El área se mide en \textbf{m$^2$}:
            «El área es 280 \textbf{m$^2$}».
            Si dices solo «metros», estás indicando una longitud.

      \vspace{0.3cm}
      \item \incorrecto\
            \textbf{«El perímetro es 4{,}800 metros cuadrados»}
            \\[6pt]
            \correcto\
            El perímetro es una longitud (1D):
            «El perímetro es 4{,}800 \textbf{m}».
            Nunca lleva el exponente 2.

      \vspace{0.3cm}
      \item \incorrecto\
            Pensar que «mayor área $\Rightarrow$ mayor perímetro».
            \\[6pt]
            \correcto\
            Son medidas \textbf{independientes}.
            Un rectángulo puede tener mayor área
            y menor perímetro que otro.
    \end{itemize}
  \end{errorbox}
\end{frame}

% ---------------------------------------------------------------
%  FRAME 13 — Resumen
% ---------------------------------------------------------------
\begin{frame}{Resumen de la sesión}
  \begin{resumenbox}
    \begin{itemize}
      \item \textbf{Perímetro} $\neq$ \textbf{Área}: miden cosas distintas.
      \item P mide el \emph{contorno} (1D, en metros).
            A mide la \emph{superficie} (2D, en m$^2$).
      \item Dos figuras con el mismo P pueden tener distintas A,
            y viceversa.
      \item Con P fijo, el cuadrado maximiza el área
            (problema isoperimérico).
      \item Rambla de Almería: $P = 840$ m, $A = 18{,}500$ m$^2$.
      \item Usar P para vallas/bordillos;
            usar A para suelos/césped.
    \end{itemize}
  \end{resumenbox}
\end{frame}

% ---------------------------------------------------------------
%  FRAME 14 — Evidencia del día
% ---------------------------------------------------------------
\begin{frame}{Evidencia del día}
  \begin{evidenciabox}
    \textbf{Ficha doble: Perímetro y Área en paralelo}

    \vspace{0.2cm}
    Elige \textbf{3 figuras} de diferente forma (rectángulo, triángulo,
    trapecio). Para cada una:
    \begin{enumerate}
      \item Dibuja la figura con sus medidas.
      \item Calcula el perímetro — muestra todos los pasos.
      \item Calcula el área — muestra todos los pasos.
      \item Escribe las unidades correctas para cada resultado.
      \item Indica para qué sirve cada medida en un contexto real.
    \end{enumerate}
    \vspace{0.2cm}
    \textbf{Criterio de éxito:} los 6 cálculos (3 P + 3 A) son correctos
    y las unidades son las adecuadas para cada tipo de medida.
  \end{evidenciabox}
\end{frame}

\end{document}



